\input{sol_pre.tex}
\input{q_2024-12-01_3.tex}
\begin{document}
\input{sol_head.tex}

\Savol{1}
$7^4=2401\equiv1\pmod{100}$, \ $2023=4\cdot505+3$: \ $7^{2023}\equiv7^3=343\equiv43\pmod{100}$. Onlar xonasi $4$.
\Javob{A}

\Savol{2}
Har bir kasr $1+\frac1n$ körinişida: \ $6-5-\left(\frac1{30}+\frac1{60}+\frac1{100}+\frac1{150}+\frac1{210}\right)$.
\[\frac{70+35+21+14+10}{2100}=\frac{150}{2100}=\frac1{14},\qquad 1-\frac1{14}=\frac{13}{14}.\]
\Javob{D}

\Savol{3}
$\text{EKUB}^2$ köpaytmaning böluvçisi: \ $\text{EKUB}\le3^2\cdot5\cdot7=315$ ($A=315$, $B=3^2\cdot5\cdot7^2\cdot11$).
\Javob{C}

\Savol{4}
$a^2+b^2=\frac{(a+b)^2+(a-b)^2}2$ — $|a-b|$ eng kiçik bölganda minimal; sonlar turli: $17$ va $19$. \ $17\cdot19=323$.
\Javob{A}

\Savol{5}
\[\frac{(2a+b)\left(4a^2-2ab+b^2\right)}{(2a-b)(2a+b)}+b=\frac{4a^2-2ab+b^2+2ab-b^2}{2a-b}=\frac{4a^2}{2a-b};\qquad \frac{4a^2}{2a-b}\cdot\frac{2a-b}{a^2}=4.\]
\Javob{C}

\Savol{6}
$\sqrt{15}+\sqrt{33}-\sqrt{22}-\sqrt{10}=\left(\sqrt{11}+\sqrt5\right)\left(\sqrt3-\sqrt2\right)$:
\[\frac{6\left(\sqrt3-\sqrt2\right)}{6\sqrt3-6\sqrt2}=1.\]
\Javob{A}

\Savol{7}
$\dfrac1{\sqrt{n+3}+\sqrt n}=\dfrac{\sqrt{n+3}-\sqrt n}3$ — yiğindi teleskopik: $\dfrac{\sqrt{a+3}-1}3=2$ $\Rightarrow$ $a=46$.
\Javob{B}

\Savol{8}
$1{,}1\cdot1{,}2=1{,}32$ — $32\%$.
\Javob{C}

\Savol{9}
$2x+3=\pm15$ $\Rightarrow$ $x=6$, $x=-9$. \ Yiğindi $-3$.
\Javob{A}

\Savol{10}
$y=\dfrac{6x+5}{x+6}$ $\Rightarrow$ $x(y-6)=5-6y$ $\Rightarrow$ $x=\dfrac{6y-5}{6-y}$. \ $f^{-1}(x)=\dfrac{6x-5}{6-x}$.
\Javob{B}

\Savol{11}
$b_{m+n}\cdot b_{m-n}=b_1^2q^{2m-2}=b_m^2$ $\Rightarrow$ $b_m=\sqrt{AB}$.
\Javob{B}

\Savol{12}
$MH=2\sqrt3$ — perpendikulyar. \ $HA=\dfrac{2\sqrt3}{\tg30^\circ}=6$, \ $HB=\dfrac{2\sqrt3}{\tg60^\circ}=2$.\\
$AB\ge HA-HB=4$ ($A$, $B$, $H$ bir nurda bölganda).
\Javob{C}

\Savol{13}
$\overrightarrow{AC_1}=\overrightarrow{AC}+\overrightarrow{CC_1}$, \ $\overrightarrow{CC_1}\perp\overrightarrow{AC}$: \ $\overrightarrow{AC}\cdot\overrightarrow{AC_1}=AC^2=2a^2=12$ $\Rightarrow$ $a=\sqrt6$.
\Javob{C}

\Savol{14}
Muntazam sakkizburçak içki burçagi $135^\circ$, oltiburçakniki $120^\circ$: \ $\alpha=135^\circ-120^\circ=15^\circ$.
\Javob{B}

\Savol{15}
Kesişiş: $(x+1)(x+2)=-(x+1)$ $\Rightarrow$ $x=-1$, $x=-3$. \ $[-3;\,-1]$ da $g\ge f$:
\[S=\int_{-3}^{-1}(-x-1)\,dx-\int_{-3}^{-1}\left|x^2+3x+2\right|dx=2-\left(\frac56+\frac16\right)=1.\]
\Javob{A}

\Savol{16}
$g(x)=\dfrac{\ln\sin x}{\ln\cos x}$, \ $g\left(\frac\pi4\right)=1<\ln\pi$ $\Rightarrow$ $\frac\pi4$ atrofida $f=\ln\pi-g$, \ $f'=-g'$.
\[g'=\frac{\ctg x\ln\cos x+\tg x\ln\sin x}{\ln^2\cos x},\qquad x=\frac\pi4:\ L=-\frac12\ln2,\ \ g'=\frac{2L}{L^2}=-\frac4{\ln2}.\]
$f'\left(\frac\pi4\right)=\dfrac4{\ln2}=4\log_2e$.
\Javob{D}

\Savol{17}
$R=4$, $BO=5$, $BC=9$. $OD\perp AB$ (uriniş): $BD=\sqrt{25-16}=3$.\\
$OD\parallel AC$ $\Rightarrow$ $\dfrac{BD}{BA}=\dfrac{OD}{AC}=\dfrac{BO}{BC}=\dfrac59$: \ $BA=\frac{27}5$, $AC=\frac{36}5$, \ $AD=\frac{12}5$.
\[DC^2=\frac{144+1296}{25}=\frac{1440}{25}\ \Rightarrow\ \frac{DC}2=\frac{6\sqrt{10}}5.\]
$O$ dan $DC$ vatargaça: $\sqrt{16-\dfrac{360}{25}}=\sqrt{\dfrac{40}{25}}=\dfrac{2\sqrt{10}}5$.
\Javob{A}

\Savol{18}
$MC=5$. \ $\triangle MNC\sim\triangle ABC$ ($\angle C$ umumiy): \ $\dfrac{MN}{AB}=\dfrac{NC}{BC}=\dfrac{MC}{AC}=\dfrac12$ $\Rightarrow$ $AC=10$.\\
$x=AC-NC=10-4=6$.
\Javob{C}

\Savol{19}
Romb tomoni $R$ va bir diagonali $R$ $\Rightarrow$ ikkita teng tomonli uçburçak: $S=\dfrac{\sqrt3}2R^2$.\\
Uçburçak: $\frac{\sqrt3}4a^2=48\sqrt3$ $\Rightarrow$ $a=8\sqrt3$, \ $R=\frac a{2\sqrt3}=4$. \ $S=\frac{\sqrt3}2\cdot16=8\sqrt3$.
\Javob{D}

\Savol{20}
$t=\log_{16}x$, \ $\log_2x-1=4t-1=4\left(t-\frac14\right)$:
\[4\left(4t^2+\frac12\right)=8\ \Rightarrow\ t^2=\frac38\ \Rightarrow\ t_1+t_2=0.\]
$x_1x_2=16^{t_1+t_2}=1$.
\Javob{D}

\Savol{21}
$a=x+1$, $b=x-2$, $a-b=3\ne0$: \ çap tomon $=\dfrac{a^{2024}-b^{2024}}{a-b}$.\\
$a^{2024}=b^{2024}$ $\Rightarrow$ $a=-b$ $\Rightarrow$ $x+1=2-x$ $\Rightarrow$ $x=\frac12$ — $1$ ta ildiz.
\Javob{D}

\Savol{22}
Aniqlaniş sohasi: $x^2-5x-6\le0$, $x\ne2$ $\Rightarrow$ $[-1;\,6]\setminus\{2\}$. Bu sohada $x+1\ge0$ — hammasi yeçim.\\
$-1+0+1+3+4+5+6=18$.
\Javob{A}

\Savol{23}
$1+2\sin x\cos x=\frac45$ $\Rightarrow$ $\sin x\cos x=-\frac1{10}$.
\[\frac{\sin^6x+\cos^6x}{\sin^2x\cos^2x}=\frac{1-3\sin^2x\cos^2x}{\sin^2x\cos^2x}=\frac{1-\frac3{100}}{\frac1{100}}=97.\]
\Javob{A}

\Savol{24}
$h-f=2$ — özgarmas $\Rightarrow$ $k=m$, $l-n=2$. \ $h\cdot f$ ning nollari $-\frac lk=-3$, $-\frac nk=-1$ (yoki aksinça), $ln=3$:\\
$l=3$, $n=1$, $k=m=1$. \ $\dfrac{k+l}{m+n}=\dfrac42=2$.
\Javob{B}

\Savol{25}
$3^x-3^{x-1}=2\cdot3^{x-1}$, \ $-2^{x+2}+2^{x-3}+2^{x-1}=-27\cdot2^{x-3}$:
\[\left(\frac32\right)^{x-1}\ge\frac{27}8\ \Rightarrow\ x\ge4.\]
\Javob{C}

\Savol{26}
$t=\sqrt[4]x\ge0$: \ $t^2-5t+6\ge0$ $\Rightarrow$ $t\le2$ yoki $t\ge3$ $\Rightarrow$ $x\in[0;\,16]\cup[81;\,\infty)$.
\Javob{C}

\Savol{27}
$\frac\pi2+3\approx4{,}57\in\left(\pi;\,\frac{3\pi}2\right)$ — III çorak.
\Javob{B}

\Savol{28}
$Ox$ ga urinadi: $R=|y_O|=5$.
\Javob{D}

\Savol{29}
$26+2\sqrt{88}=\left(\sqrt{22}+2\right)^2$: \ $\dfrac{\left(\sqrt{22}-2\right)\left(\sqrt{22}+2\right)}9=\dfrac{18}9=2$.
\Javob{C}

\Savol{30}
$a_n=1+(2+4+\dots+2(n-1))=1+n(n-1)=4033$ $\Rightarrow$ $n(n-1)=4032=64\cdot63$ $\Rightarrow$ $n=64$.
\Javob{D}

\Savol{31}
Jami: $5\cdot5\cdot4=100$. Toq: oxiri $1$ yoki $9$ ($2$), birinçi raqam ($0$ va oxirgisidan boşqa) $4$, örtasi $4$: $32$.
\[P=\frac{32}{100}=\frac8{25}.\]
\Javob{B}

\Savol{32}
$A\cap B=\varnothing$ $\Rightarrow$ $(A\cup B)'\cup A=B'$. \ $n(U)=10$, $n(B)=3$: \ $2^7=128$.
\Javob{C}

\newpage
\Savol{33}
Qirqilgan qism ($x$, $y$ — çap pastki burçakdan, cm): $[0;3]\times[0;9]$ töğri törtburçak, $3\le x\le7$ da radiusi $4$ çorak doira «taşqarisi», $7\le x\le12$ da radiusi $5$ çorak doira:
\[S_{\text{qirq}}=3\cdot9+\left(4\cdot9-\frac{3\cdot16}4\right)+\frac{3\cdot25}4=27+24+18{,}75=69{,}75.\]
\textbf{33.} $S=225-69{,}75=155{,}25$ cm$^2$ — \textbf{A}.

\textbf{34.} Kesim perimetri: $15+15+3+6+3+\frac{2\cdot3\cdot4}4+\frac{2\cdot3\cdot5}4=55{,}5$ cm.\\
$S_{\text{töla}}=2\cdot155{,}25+55{,}5\cdot100=5860{,}5$ cm$^2=58{,}605$ dm$^2$ — \textbf{D}.

\textbf{35.} $V=155{,}25\cdot100=15525$ cm$^3$ — \textbf{B}.
\Javob{33 — A, \ 34 — D, \ 35 — B}

\newpage
\Savol{36}
$u=x^2+k$: \ $(u-2x)^2+3x(u+2x)=30x^2$ $\Rightarrow$ $u^2-ux-20x^2=0$ $\Rightarrow$ $(u-5x)(u+4x)=0$:
\[x^2-5x+k=0\ \ (D_1=25-4k),\qquad x^2+4x+k=0\ \ (D_2=16-4k).\]
Umumiy ildiz faqat $x=0$ ($k=0$).

$k>\frac{25}4$: ildiz yöq; \ $k=\frac{25}4$: $1$ ta; \ $4<k<\frac{25}4$: $2$ ta; \ $k=4$: $3$ ta ($1;4;-2$); \ $k<4$: $4$ ta, $k=0$ da $3$ ta ($0;5;-4$).

a) $k>6{,}25$ $\Rightarrow$ eng kiçik butun $k=7$.

b) $k_1=\frac{25}4$, \ $k_2=6$, \ $k_3=4$, \ $k_4=3$:
\[k_1k_3(k_2+k_4)=\frac{25}4\cdot4\cdot9=225.\]
\Javob{a) 7; \quad b) 225}

\newpage
\Savol{37}
Aniqlaniş sohasi: $x^2-6x+5\le0$ $\Rightarrow$ $x\in[1;\,5]$.

1) Ildiz ostidagi ifoda $0$: $x=1$, $x=5$.

2) $\cos2x=\frac12$: $2x=\pm\frac\pi3+2\pi n$, \ $2x\in[2;\,10]$: \ $2x=\frac{5\pi}3$, $2x=\frac{7\pi}3$ $\Rightarrow$ $x=\frac{5\pi}6\approx2{,}62$, $x=\frac{7\pi}6\approx3{,}67$.

a) $4$ ta. \quad b) Eng kiçik musbat ildiz $1$.
\Javob{a) 4; \quad b) 1}

\newpage
\Savol{38}
$BE=x$: \ $AE=\sqrt{900+x^2}$, \ $EC=80-x$.
\[T(x)=\frac{\sqrt{900+x^2}}3+\frac{80-x}5,\qquad T'=\frac x{3\sqrt{900+x^2}}-\frac15=0\ \Rightarrow\ 5x=3\sqrt{900+x^2}\ \Rightarrow\ x=22{,}5.\]
$AE=\sqrt{900+506{,}25}=37{,}5$, \ $EC=57{,}5$.

a) $T=\dfrac{37{,}5}3+\dfrac{57{,}5}5=12{,}5+11{,}5=24$ soat.

b) $EC=57{,}5$ m.
\Javob{a) 24; \quad b) $57{,}5$}

\newpage
\Savol{39}
$AH=3k$, $HB=2k$ $\Rightarrow$ $AD=AB=5k$, \ $DH=4k$. Doira diametri — romb balandligi: $r=2k$.

a) $\dfrac{S_{\text{romb}}}{S_{\text{doira}}}=\dfrac{5k\cdot4k}{3\cdot4k^2}=\dfrac53$.

b) $O$ — diagonallar kesişgan nuqta, $\angle DOC=90^\circ$, $OM\perp DC$: \ $DM=\dfrac{OD^2}{DC}$, $CM=\dfrac{OC^2}{DC}$ $\Rightarrow$ $\dfrac{CM}{MD}=\dfrac{AC^2}{BD^2}$.
\[BD^2=DH^2+HB^2=20k^2,\qquad AC^2=4AB^2-BD^2=80k^2\ \Rightarrow\ \frac{CM}{MD}=4.\]
\Javob{a) $\frac53$; \quad b) 4}

\newpage
\Savol{40}
Silindr radiusi $=$ oltiburçak tomoni $a$, balandliklar $h$. Piramida apofemasi $l$, apofema proyeksiyasi $\frac{a\sqrt3}2$.
\[\frac{2\pi ah}{6\cdot\frac12a\cdot l}=\frac{2\pi}3\cdot\frac hl=\frac{\sqrt2\pi}3\ \Rightarrow\ \frac hl=\frac{\sqrt2}2.\]
a) $V_{\text{sil}}=\pi a^2h$, \ $V_{\text{pir}}=\frac13\cdot\frac{3\sqrt3}2a^2h=\frac{\sqrt3}2a^2h$: \ $\dfrac{2\pi}{\sqrt3}$.

b) $\sin\varphi=\dfrac hl=\dfrac{\sqrt2}2$ $\Rightarrow$ $\varphi=45^\circ$.
\Javob{a) $\frac{2\pi}{\sqrt3}$; \quad b) $45^\circ$}

\newpage
\Savol{41}
$x\to\frac1x$: \ $f\left(\frac1x\right)+4f(x)=x$. \ $4$ ga köpaytirib birinçisini ayiramiz:
\[15f(x)=4x-\frac1x\ \Rightarrow\ f(x)=\frac1{15}\left(4x-\frac1x\right).\]
a) $f\left(\frac12\right)=\frac1{15}(2-2)=0$.

b) $f'(x)=\frac1{15}\left(4+\frac1{x^2}\right)$, \ $f'(1)=\dfrac5{15}=\dfrac13$.
\Javob{a) 0; \quad b) $\frac13$}

\newpage
\Savol{42}
$\angle B=2\alpha$: \ $AB=8\cos2\alpha$, \ $BL=\dfrac{AB}{\cos\alpha}=\dfrac{8\cos2\alpha}{\cos\alpha}=\dfrac8{\sqrt3}$:
\[\sqrt3\left(2\cos^2\alpha-1\right)=\cos\alpha\ \Rightarrow\ 2\sqrt3c^2-c-\sqrt3=0\ \Rightarrow\ c=\frac{\sqrt3}2,\ \alpha=30^\circ.\]
$\angle B=60^\circ$: \ $AB=4$, \ $AC=4\sqrt3$.

a) Kiçik katet $AB=4$.

b) $AM=2$: \ $CM=\sqrt{4+48}=2\sqrt{13}$.
\Javob{a) 4; \quad b) $2\sqrt{13}$}

\newpage
\Savol{43}
Aylana radiusi $r$ — sakkizburçakka içki, oltiburçakka taşqi: $b=r$.
\fig{\begin{tikzpicture}[scale=1.1,font=\small]
\draw (22.5:2.165)--(67.5:2.165)--(112.5:2.165)--(157.5:2.165)--(202.5:2.165)--(247.5:2.165)--(292.5:2.165)--(337.5:2.165)--cycle;
\draw (0,0) circle (2); \draw (0:2)--(60:2)--(120:2)--(180:2)--(240:2)--(300:2)--cycle;
\draw[axs] (0,0)--(0,2); \draw[ax] (0,0)--(67.5:2.165); \fill (0,0) circle (1.5pt);
\node[below] at (0,0) {$O$}; \node[aux,left] at (0,1.2) {$r$}; \node[aux] at (0.62,2.25) {$\frac a2$};
\end{tikzpicture}}
a) $\dfrac a2=r\tg22{,}5^\circ=r\left(\sqrt2-1\right)$ $\Rightarrow$ $\dfrac ab=2\left(\sqrt2-1\right)=2\sqrt2-2$.

b) $\dfrac{3\sqrt3}2b^2=3\sqrt6+3\sqrt3$ $\Rightarrow$ $b^2=2\left(\sqrt2+1\right)$. \ $a^2=4\left(\sqrt2-1\right)^2b^2=8\left(\sqrt2-1\right)$.
\[S_8=2\left(1+\sqrt2\right)a^2=16\left(\sqrt2+1\right)\left(\sqrt2-1\right)=16.\]
\Javob{a) $2\sqrt2-2$; \quad b) 16}

\newpage
\Savol{44}
$(y-2)^2+x^2=4$, $y\ge2$ — markazi $Q(0;2)$, radiusi $2$ bölgan yuqori yarim aylana.
\fig{\begin{tikzpicture}[scale=0.8,font=\small]
\draw[-{Stealth[length=4pt]}] (-2.8,0)--(2.8,0) node[right] {$x$}; \draw[-{Stealth[length=4pt]}] (0,-0.4)--(0,4.6) node[left] {$y$};
\draw[thick] (2,2) arc (0:180:2); \draw[ax] (0,0)--(2,2); \draw[ax] (0,0)--(-2,2); \fill (0,2) circle (1.5pt);
\node[below left] at (0,0) {$O$}; \node[right] at (0,2) {$Q$}; \node[right] at (2,2) {$(2;2)$};
\end{tikzpicture}}
a) $O$ nuqta aylananing pastki qismida yotadi ($OQ=2=R$); yuqori yarim aylana nuqtalaridan eng yaqini uçlari $(\pm2;\,2)$: \ $d=\sqrt{4+4}=2\sqrt2$.

b) $\displaystyle\int_0^2\sqrt{4-x^2}\,dx$ — çorak doira yuzi $\pi$: \ $\pi+2\cdot2=\pi+4$.
\Javob{a) $2\sqrt2$; \quad b) $\pi+4$}

\newpage
\Savol{45}
Parabola: asosi $6$, balandligi $4$.

a) Parabolik segment yuzi: $\dfrac23\cdot6\cdot4=16$.

b) Örtadan $t$ masofada parabola: $y=4-\frac49t^2$. $A$ dan örtagaça $5$. $AC$ ning oğmasi $m$: \ $m(5-t)=4-\frac49t^2$ uriniş şarti — diskriminant nol:
\[\frac49t^2-mt+5m-4=0:\ \ m^2-\frac{16}9(5m-4)=0\ \Rightarrow\ 9m^2-80m+64=0\ \Rightarrow\ m=\frac89\ \ (m=8\ \text{da uriniş nuqtasi parabola yoyidan taşqarida}).\]
$C$ ning balandligi $5\cdot\frac89=\frac{40}9$: \ $S_{ABC}=\frac12\cdot10\cdot\frac{40}9=\dfrac{200}9=22\dfrac29$.
\Javob{a) 16; \quad b) $22\frac29$}

\ShowJavoblar
\end{document}
